% ==============================================================================
% Springer Nature journal submission template (skeleton)
% Official class: sn-jnl.cls  (download from Springer Nature LaTeX template
% page: https://www.springernature.com/gp/authors/campaigns/latex-author-help
% or search "Springer Nature LaTeX Template" on Overleaf). The sn-jnl.cls is
% NOT distributed in standard TeX Live, so to compile this skeleton out-of-the
% box we use the standard `article` class with Springer-editorial styling. For
% an actual Springer submission, replace the preamble below with:
%
%   \documentclass[sn-basic]{sn-jnl}      % numbered (Basic)
%   \documentclass[sn-mathphys]{sn-jnl}   % math/physics numbered
%   \documentclass[sn-apa]{sn-jnl}        % APA author-year
%   \documentclass[sn-standardnature]{sn-jnl}  % Nature in-house style
%   \documentclass[sn-vancouver]{sn-jnl}  % Vancouver numbered
%   \documentclass[default,iicol]{sn-jnl} % two-column review style
%
% Editorial style (Springer Nature, "Key Style Points"):
%   Layout     : single column, 10pt, double-spaced for review
%   Body font  : Times Roman
%   Math font  : Times (mathptmx)
%   Citation   : numbered [1] OR author-year (Author 20YY) per journal
%   Figures    : single column ~85 mm  |  double column ~170 mm
%   Captions   : below figures, above tables
% References  : sn-basic.bst, sn-mathphys.bst, sn-apa.bst, sn-vancouver.bst
%               (shipped with the sn-jnl template archive)
% ==============================================================================
\documentclass[11pt,a4paper]{article}

% ---- To use the real Springer class, comment out the line above and use: ----
% \documentclass[sn-basic]{sn-jnl}
% \usepackage{sn-jnl.cls}  % (the cls is loaded by \documentclass, not \usepackage)

\usepackage[a4paper,top=30mm,bottom=30mm,left=30mm,right=30mm]{geometry}
\usepackage[T1]{fontenc}
\usepackage[utf8]{inputenc}
\usepackage{mathptmx}            % Times text + math
\usepackage{helvet}
\usepackage{courier}
\usepackage{amsmath,amssymb}
\usepackage{graphicx}
\usepackage{booktabs}
\usepackage{array}
\usepackage{xcolor}
\usepackage{lineno}

% ---- Citations: natbib (Springer uses natbib for both numbered and author-year) ----
\usepackage[numbers,sort&compress]{natbib}   % numbered style (sn-basic)
% For author-year (sn-apa) use instead:
% \usepackage[authoryear]{natbib}
\bibliographystyle{unsrtnat}                 % numeric, citation order
% With the real sn-jnl class, use:
%   \bibliographystyle{sn-basic}      % numbered
%   \bibliographystyle{sn-mathphys}   % math/phys numbered
%   \bibliographystyle{sn-apa}        % APA author-year
%   \bibliographystyle{sn-vancouver}  % Vancouver

\linenumbers                     % review mode

% \usepackage{hyperref}           % load last if used

\title{A Concise and Informative Title for the Springer Nature Submission}
\author{First Author\textsuperscript{1,*} \and Second Author\textsuperscript{2} \and Senior Author\textsuperscript{1}}
\date{}

\begin{document}
\maketitle

% ---- Springer author block (mirrors sn-jnl \author*\affiliation* layout) ----
\noindent\textsuperscript{1}Department of X, University of Y, City, Country\\
\textsuperscript{2}Department of Z, Institute of W, City, Country\\
\textsuperscript{*}Corresponding author: first.author@univ.edu

\vspace{1em}

% ==============================================================================
% Abstract  (Springer: structured or unstructured, ~150-250 words, no cites)
% ==============================================================================
\begin{abstract}
% TODO: 150-250 words. Many Springer journals request a structured abstract
% (Background / Methods / Results / Conclusions); check the journal's guide
% for authors. No citations, no undefined abbreviations.
\end{abstract}

\noindent\textbf{Keywords:} % TODO: 4-6 keywords separated by commas.
keyword one, keyword two, keyword three

\vspace{1em}

% ==============================================================================
% 1. Introduction
% ==============================================================================
\section{Introduction}
\label{sec:intro}
% TODO: context -> gap -> approach -> contributions -> roadmap. Use \citep{}
% (parenthetical) and \citet{} (textual) for author-year, or \cite{} for numbered.

% ==============================================================================
% 2. Related Work
% ==============================================================================
\section{Related Work}
\label{sec:related}
% TODO: cluster prior work and position your contribution.

% ==============================================================================
% 3. Methods
% ==============================================================================
\section{Methods}
\label{sec:methods}
% TODO: problem formulation, proposed method, theoretical justification.

\subsection{Problem formulation}
% TODO.

\subsection{Proposed method}
% TODO.

% ==============================================================================
% 4. Experiments
% ==============================================================================
\section{Experiments}
\label{sec:experiments}
% TODO: setup, main results, ablation, qualitative examples, statistics.

\begin{figure}[t]
\centering
% \includegraphics[width=0.9\textwidth]{figures/result.pdf}
\fbox{\rule{0pt}{5cm}\rule{0.7\textwidth}{0pt}}
\caption{Figure caption goes BELOW the figure (Springer convention).
Single-column figures ~85 mm wide, double-column figures ~170 mm wide.}
\label{fig:result}
\end{figure}

\begin{table}[t]
\caption{Table caption goes ABOVE the table (Springer convention).}
\label{tab:results}
\centering
\begin{tabular}{lcc}
\toprule
Method   & Metric A       & Metric B       \\
\midrule
Baseline & 0.812          & 0.754          \\
Ours     & \textbf{0.891} & \textbf{0.832} \\
\bottomrule
\end{tabular}
\end{table}

% ==============================================================================
% 5. Conclusion
% ==============================================================================
\section{Conclusion}
\label{sec:conclusion}
% TODO: recap contribution and headline result; one sentence of future work.

% ==============================================================================
% Back matter
% ==============================================================================
\section*{Acknowledgements}
% TODO: funding, grant numbers.

\section*{Funding}
% TODO: grant IDs and agencies (Springer requires an explicit Funding section).

\section*{Competing Interests}
% TODO: declare competing interests, or "The authors have no competing interests."

\section*{Data Availability}
% TODO: state where data and code are deposited (DOIs preferred).

\section*{Author Contributions}
% TODO: e.g. "F.A. conceived the study, performed the analysis, and wrote the
% manuscript. S.A. supervised the project."

% ==============================================================================
% References
% Compile:  pdflatex main  ->  bibtex main  ->  pdflatex main  ->  pdflatex main
% ==============================================================================
\bibliography{references}
% Inline fallback (uncomment if not using BibTeX):
% \begin{thebibliography}{99}
% \bibitem{ref_example} Author, F.: Title of the paper. J. Name XX, 1--10 (20YY)
% \end{thebibliography}

\end{document}
